\section{Positron Emission Tomography (PET)}

\subsection{How PET works}

\subsection{Deterministic reconstruction methods}

\paragraph{Radon transform}

\begin{definition}{Radon Transform}
\label{def:radon_transform}
    The Radon transform is an integral transform that takes a two-dimensional function $f(x,y)$ and computes its line integrals along straight lines parameterized by angle $\phi$ and distance $s$ from the origin.
    Mathematically, the Radon transform is defined as:
    \begin{equation}
        g(s, \phi) = \int \int_{-\infty}^{\infty} f(x,y) \, \delta(s - x \cos \phi - y \sin \phi) \, dx dy
    \end{equation}
\end{definition}

In the context of PET imaging, the Radon transform models the process of acquiring projection data (i.e., sinograms) from the distribution of radiotracer within the body.

\begin{theorem}{Central Slice Theorem}
\label{thm:central_slice}
    The Radon transform is invertible in the frequency domain.
\end{theorem}
\begin{proof}
    We consider the 1D Fourier transform of the projection $g(s,\phi)$ with respect to the variable $s$:
    \begin{equation*}
        G(\omega, \phi) = \int_{-\infty}^{\infty} g(s, \phi) e^{-i \omega s} ds
    \end{equation*}
    Substituting the expression for $g(s,\phi)$ into the Fourier transform, we get:
    \begin{equation*}
        G(\omega, \phi) = \int_{-\infty}^{\infty} \left( \int \int_{-\infty}^{\infty} f(x,y) \, \delta(s - x \cos \phi - y \sin \phi) \, dx dy \right) e^{-i \omega s} ds
    \end{equation*}
    Interchanging the order of integration, we have:
    \begin{equation*}
        G(\omega, \phi) = \int \int_{-\infty}^{\infty} f(x,y) \left( \int_{-\infty}^{\infty} \delta(s - x \cos \phi - y \sin \phi) e^{-i \omega s} ds \right) dx dy
    \end{equation*}
    The inner integral evaluates to $e^{-i \omega (x \cos \phi + y \sin \phi)}$, leading to:
    \begin{equation*}
        G(\omega, \phi) = \int \int_{-\infty}^{\infty} f(x,y) e^{-i \omega (x \cos \phi + y \sin \phi)} dx dy
    \end{equation*}
    This expression represents the 2D Fourier transform of $f(x,y)$ evaluated at the frequency coordinates $(\omega \cos \phi, \omega \sin \phi)$. Thus, we have:
    \begin{equation}
        G(\omega, \phi) = F(\omega \cos \phi, \omega \sin \phi)
    \end{equation}
    where $F(u,v)$ is the 2D Fourier transform of $f(x,y)$. This shows that the 1D Fourier transform of the projection at angle $\phi$ corresponds to a slice of the 2D Fourier transform of the object $f(x,y)$ along the direction defined by $\phi$.
\end{proof}


This theorem comes with limitations in practice as we can only acquire a limited number of projections and the data is noisy. Indeed, going from the frequency domain to the spatial domain requires interpolation, which can introduce artifacts if the sampling is not dense enough.
Hence, it is not possible to perfectly reconstruct the original image using only the Radon transform and its inverse. 

\paragraph{Filtered Back Projection (FBP)}
Another possible reconstruction method is known as Back Projection (BP).
This method consists in smearing back each projection across the image along the direction it was acquired.

\begin{definition}{Back Projection}
    \label{def:back_projection}
    The Back projection operation $b$ is defined as follows:
    \begin{equation}
        b(x, y) = \int_{0}^{\pi} g(s, \phi) d\phi = \int_{0}^{\pi} g(x \cos \phi + y \sin \phi, \phi) d\phi
    \end{equation}
\end{definition}

\begin{proposition}
    The Back Projection operation can be interpreted as a convolution of the original image with a $1/r$ kernel, where $r$ is the distance from the center of the image. This convolution results in a blurring effect, as high-frequency components of the image are attenuated
\end{proposition}

\begin{proof}
    Taking the 2D Fourier transform of the Back Projection operation, we have:
    \begin{align*}
        B(u, v) &= \int_{-\infty}^{\infty} \int_{-\infty}^{\infty} \int_{0}^{\pi} g(s, \phi) e^{-2i \pi (ux + vy)} dx dy d\phi\\
        &= \int_{0}^{\pi} \left( \int_{-\infty}^{\infty} \int_{-\infty}^{\infty} g(s, \phi) e^{-2i \pi (ux + vy)} dx dy \right) d\phi\\
        &= \int_{0}^{\pi} \left( \int_{-\infty}^{\infty} g(s, \phi) e^{-2i \pi s \sqrt{u^2 + v^2}} ds \right) d\phi\\
        &= \int_{0}^{\pi} G(\sqrt{u^2 + v^2}, \phi) d\phi
    \end{align*}
    Substituting the expression for $G(\omega, \phi)$ from the Central Slice Theorem (\ref{thm:central_slice}), we get:
    \begin{equation*}
        B(u,v) = \int_{0}^{\pi} F(\sqrt{u^2 + v^2} \cos \phi, \sqrt{u^2 + v^2} \sin \phi) d\phi
    \end{equation*}
    This integral represents a radial integration of the 2D Fourier transform of the original image $f(x,y)$.
    The result is that the Back Projection operation effectively applies a low-pass filter to the original image, attenuating high-frequency components and resulting in a blurred image.
\end{proof}

To counteract this blurring effect, a filtering step is introduced before the Back Projection, leading to the Filtered Back Projection (FBP) method.

\begin{definition}{Filtered Back Projection}
    The Filtered Back Projection operation consists in applying a high-pass filter to the projections before performing the Back Projection:
    \begin{equation}
        f_{\text{FBP}}(x, y) = \int_{0}^{\pi} \left( g(s, \phi) * h(s) \right) d\phi
    \end{equation}
    where $h(s)$ is a high-pass filter kernel, such as the Ram-Lak filter.
\end{definition}

\subsection{Iterative reconstruction in PET}

\subsubsection{Poisson bayesian inference}

In PET imaging, the data acquisition process can be modeled as a Poisson process as the number of detected events in each detector pair follows a Poisson distribution.
Thus, the likelihood of observing the measured data $y$ given the image $x$ can be expressed as:
\begin{equation}
    p(y|x) = \prod_{i=1}^{d} \frac{(y_i(x))^{y_i} e^{-y_i(x)}}{y_i!}
\end{equation}
where $y_i(x) = (Ax)_i$ is the expected number of counts in detector pair $i$ given the image $x$. It is modeled as a linear operation of the image through the system matrix $A$.

We want to estimate the image $x$ that maximizes the likelihood of the unobserved data $x$ given the observed data $y$. This is known as the Maximum Likelihood (ML) estimation. We will transform this formulation into a minimization problem by taking the negative log-likelihood:

\begin{equation}
    \hat{x}_{\text{ML}} = \arg \min_{x} \sum_{i=1}^{d} \left( (Ax)_i - y_i \log (Ax)_i\right)
    \label{eq:ml_estimation}
\end{equation}

\subsubsection{Maximum-Likelihood Expectation-Maximization (MLEM)}

Among iterative algorithms for PET reconstrcution, the \textbf{Maximum-Likelihood Expectation-Maximization (MLEM)} proposed by \cite{Shepp1982} is one of the most widely used. It is an iterative method to find the ML estimate of the image $x$ given the measured data $y$.

The negative log-likelihood is convex, thus the MLEM algorithm is guaranteed to converge to a global minimum, providing an estimate of the image $x$ that is consistent with the measured data $y$ under the Poisson noise model.
However, it can be slow to converge and can produce noisy images if the number of iterations is too high.
In order to address this, iterations can be stopped early, or output can be post-processed or post-filtered.

\begin{algorithm}
\caption{Maximum-Likelihood Expectation-Maximization (MLEM)}
\begin{algorithmic}[1]
\State \textbf{Input:} Measured data $y$, system matrix $A$, initial estimate $x^{(0)}$
\State \textbf{Output:} Estimated image $x^{(k)}$
\For{$k = 0$ to $K-1$}
    \State \textbf{E-step:} Compute the expected number of counts
    \begin{equation*}
        \hat{y}_i^{(k)} = (Ax^{(k)})_i
    \end{equation*}
    \State \textbf{M-step:} Update the image estimate
    \begin{equation*}
        x_j^{(k+1)} = x_j^{(k)} \frac{\sum_{i=1}^{d} A_{ij} \frac{y_i}{\hat{y}_i^{(k)}}}{\sum_{i=1}^{d} A_{ij}}
    \end{equation*}
\EndFor
\State \textbf{Return} $x^{(K)}$
\end{algorithmic}
\label{alg:mlem}
\end{algorithm}

% \begin{proposition}
%     The MLEM algorithm converges to a local maximum of the likelihood function, providing an estimate of the image $x$ that is consistent with the measured data $y$ under the Poisson noise model.
% \end{proposition}

% \begin{proof}
%     First we compute the complete data ; $z_{ij}$ is the number of photons emitted from voxel $j$ and detected by detector pair $i$. Hence, $y_i$ = $\sum_{j} z_{ij}$ and $\bar{z_{ij}} = a_{ij} x_j$.
%     The complete data likelihood is given by:
%     \begin{align*}
%         &zp(z|x) = \prod_{i=1}^{d} \prod_{j=1}^{n} \frac{(a_{ij} x_j)^{z_{ij}}}{z_{ij}!} e^{-a_{ij} x_j} \\
%         & \Rightarrow \log p(z|x) = \sum_{i=1}^{d} \sum_{j=1}^{n} \left( z_{ij} \log (a_{ij} x_j) - a_{ij} x_j + K \right) 
%     \end{align*}
%     Then we compute its derivative with respect to $x_j$ and set it to zero to find the maximum:
%     \begin{equation*}
%         \frac{\partial \log p(z|x)}{\partial x_j} = \sum_{i=1}^{d} \left( \frac{z_{ij}}{x_j} - a_{ij} \right) = 0
%     \end{equation*}
%     This leads to the update rule:
%     \begin{equation*}
%         x_j = \frac{\sum_{i=1}^{d} z_{ij}}{\sum_{i=1}^{d} a_{ij}}
%     \end{equation*}
%     Since $z_{ij}$ is unobserved, we replace it with its expectation given the current estimate $x^{(k)}$:
%     \begin{equation*}
%         \hat{z}_{ij}^{(k)} = \mathbb{E}[z_{ij} | y, x^{(k)}] = a_{ij} x_j^{(k)} \frac{y_i}{(Ax^{(k)})_i}
%     \end{equation*}
%     Substituting this into the update rule gives the M-step of the MLEM algorithm:
%     \begin{equation*}
%         x_j^{(k+1)} = x_j^{(k)} \frac{\sum_{i=1}^{d} a_{ij} \frac{y_i}{(Ax^{(k)})_i}}{\sum_{i=1}^{d} a_{ij}}
%     \end{equation*}
%     This iterative process continues until convergence, leading to an estimate of $x$ that maximizes the likelihood function.
% \end{proof}

\subsubsection{Ordered Subset Expectation Maximization (OSEM)}

In order to speed up the convergence of the MLEM algoritm,
the \textbf{Ordered Subset Expectation Maximization (OSEM)} algorithm was proposed by \cite{Hudson1994}. It is a variant of the MLEM algorithm that divides the measured sinogram $y$ into subsets $\left\{y_1, \hdots, y_Q\right\}$ and performs the E-step and M-step from \ref{alg:mlem} on each subset sequentially. Algorithm is given in \ref{alg:osem}.

As the algorithm is not designed to minimized a well-defined cost funcion, the latter does not lead to the
same ML estimate of \ref{eq:ml_estimation} as the MLEM algorithm. However, it is much faster to converge and can be used in practice with a large number of subsets without significant loss of image quality.
Post-processing or post-filtering can also be applied to the output of the OSEM algorithm to further improve image quality.

\begin{algorithm}
\caption{Ordered Subset Expectation Maximization (OSEM)}
\begin{algorithmic}[1]
\State \textbf{Input:} Measured data $y$, system matrix $A$, initial estimate $x^{(0)}$, number of subsets $Q$
\State \textbf{Output:} Estimated image $x^{(k)}$
\For{$k = 0$ to $K-1$}
    \For{$q = 1$ to $Q$}
        \State \textbf{E-step:} Compute the expected number of counts for subset $q$
        \begin{equation*}
            \hat{y}_{i}^{(k,q)} = (Ax^{(k)})_i \quad \text{for } i \in \text{subset } q
        \end{equation*}
        \State \textbf{M-step:} Update the image estimate using subset $q$
        \begin{equation*}
            x_j^{(k,q)} = x_j^{(k,q-1)} \frac{\sum_{i \in \text{subset } q} A_{ij} \frac
            {y_i}{\hat{y}_i^{(k,q)}}}{\sum_{i \in \text{subset } q} A_{ij}}
        \end{equation*}
    \EndFor
\EndFor
\State \textbf{Return} $x^{(K)}$
\end{algorithmic}
\label{alg:osem}
\end{algorithm}

\subsubsection{Regularized reconstruction in PET}

This reconstruction problem can also be formulatted as a regularized optimization by introducing some prior knowledge on the image $x$.
Indeed, Bayes's rule gives us the following expression for the posterior distribution of $x$ given $y$:
\begin{equation}
    p(x|y) = \frac{p(y|x) p(x)}{p(y)}
\end{equation}
where $p(x)$ is the prior distribution of the image $x$ and $p(y)$ is the marginal likelihood of the data $y$.
Once agin, by taking the negative log-likelihood of $p(x|y)$, we can transform the problem into a penalized likelihood minimization problem:
\begin{equation}
    \hat{x}_{\text{MAP}} = \arg \min_{x} \left( -\log p(y|x) - \log p(x) \right)
\end{equation}
where $-\log p(y|x)$ is the data fidelity term and $-\log p(x)$ is the regularization term that encodes our prior knowledge on the image $x$.
This prior can be expressed in the following exponential form :

\begin{equation}
    p(x) \propto e^{-\beta R(x)}
\end{equation}
where $R(x)$ is a regularization function and $\beta$ is a regularization parameter that controls the trade-off between data fidelity and regularization. The optimisation hence becomes:

\begin{equation}
    \hat{x}_{\text{MAP}} = \arg \min_{x} \left( -\log p(y|x) + \beta R(x) \right)
\end{equation}

The penalized likelihood approach allows us to incorporate various types of prior knowledge about the image, such as smoothness, sparsity, or anatomical information.
The choice of the regularization function $R(x)$ must alleviate the ill-posedness of the reconstruction problem while preserving important features of the image, such as edges and textures.
The following paragraph presents some commonly used regularization functions in PET image reconstruction.

\paragraph{Markov Random Fields}

\textbf{Markov Random Fields (MRFs)} are a class of probabilistic graphical models that can be used to model spatial dependencies in images. In the context of PET image reconstruction, MRFs can be used to define a prior distribution over the image $x$ that encourages spatial smoothness while preserving edges.
Its most basinc formulation can be expressed as follows:

\begin{equation}
    R(x) = \sum_{j} \sum_{k \in N(j)} w_{jk} \phi(x_j, x_k)
    \label{eq:mrf}
\end{equation}
where $N(j)$ is the set of neighboring pixels of pixel $j$, $w_{jk}$ is a weight that controls the strength of the interaction between pixels $j$ and $k$, and $\phi$ is a potential function that defines the type of regularization applied to the differences between neighboring pixels.

\subparagraph{$l_2$ penalty (Quadratic)}

One of the most basic choice for the potential function $\phi$ is the quadratic function, first proposed by \cite{tikhonov1977solutions}, known as \textbf{Tikhonov regularization}. It encourages global smoothness by penalizing the squared differences between neighboring pixels. The side effect of this regularization is that it can oversmooth the image and blur edges, which are important features in medical imaging as we want to distinguish between different tissues and structures.
\begin{equation}
    \phi(x_j, x_k) = (x_j - x_k)^2
\end{equation}

\subparagraph{$l_1$ penalty (Total Variation)}

In order to preserve edges while still encouraging smoothness, the \textbf{Total Variation (TV)} regularization was proposed by \cite{Rudin1992}. It penalizes the absolute differences between neighboring pixels, which allows for sharp transitions in the image while still reducing noise in homogeneous regions.
\begin{equation}
    \phi(x_j, x_k) = |x_j - x_k|
\end{equation}

\subparagraph{Huber penalty}

Hybrid penalties combine the advantages of both $l_2$ and $l_1$ regularization by applying a quadratic penalty to small differences and a linear penalty to large differences. This allows for smoothness in homogeneous regions while preserving edges. One example of a hybrid penalty is the Huber function, which is defined as follows:
\begin{equation}
    \phi(x_j, x_k) = \begin{cases}
        \displaystyle \frac{(x_j - x_k)^2}{2 \delta}  & \text{if } |x_j - x_k| \leq \delta \\
        \displaystyle |x_j - x_k| - \frac{\delta}{2} & \text{if } |x_j - x_k| > \delta
    \end{cases}
\end{equation}

\subparagraph{Gibbs penalty}

The Gibbs penalty is a concave prior that has been shown to be effective in PET imaging \cite{Nuyts2002}. It is also known as the \textbf{Relative Differences}.

\begin{equation}
    \phi(x_j, x_k) = \frac{(x_j - x_k)^2}{x_j + x_k + \gamma |x_j - x_k|}
\end{equation}
where the $\gamma$ parameter controls which relative diferences are considered as large.

\subsection{Deep learning for PET image reconstruction}

\subsubsection{Supervised methods}
\subsubsection{Unrolling techniques}